\documentclass{article}

% Language setting
% Replace `english' with e.g. `spanish' to change the document language

\usepackage[letterpaper,top=2cm,bottom=5cm,left=3cm,right=3cm,marginparwidth=1.75cm]{geometry}
\usepackage{fontawesome5}
% Useful packages
\usepackage{amsmath}
\usepackage{graphicx}
\usepackage[colorlinks=true, allcolors=blue]{hyperref}
\usepackage{amssymb,stackengine}
\usepackage{fancyhdr}
\usepackage{lipsum} % Para generar texto de ejemplo
\usepackage{multicol}
\pagestyle{fancy}
% PAQUETES QUE NO SON DEL FORMATO
\usepackage{tabularx}
\usepackage{booktabs}
\usepackage{float}
\usepackage{adjustbox}
\usepackage{graphicx}   % Para imagenes
\usepackage{caption}    % Para captionof sin figure flotante
\usepackage{listings}
\usepackage{xcolor}
\usepackage[spanish]{babel} % Si quieres que "Figura" salga en espanol

\lstdefinestyle{pythonstyle}{
    language=Python,
    basicstyle=\ttfamily\small,
    breaklines=true,
    breakatwhitespace=false,
    columns=fullflexible,
    keepspaces=true,
    showstringspaces=false,
    frame=single,
    framerule=0.3pt,
    tabsize=4,
    xleftmargin=2mm,
    xrightmargin=2mm,
    keywordstyle=\color{blue!60!black},
    commentstyle=\color{green!40!black},
    stringstyle=\color{red!50!black}
}
\lstset{style=pythonstyle}

\newcommand{\checkedbox}{\stackinset{c}{0pt}{c}{0pt}{\checkmark}{$\square$}}
\newcommand{\espacioimagen}[2]{%
\IfFileExists{#1}{%
\includegraphics[width=0.92\textwidth]{#1}%
}{%
\fbox{\parbox[c][5.8cm][c]{0.90\textwidth}{\centering Espacio reservado para #2\\(agregar imagen en #1)}}%
}%
}

\fancyhf{} % Limpiar encabezado y pie de pagina por defecto
\setlength{\headheight}{70pt} % Ajuste del tamano del encabezado (puedes aumentar este valor si es necesario)

% Definicion del encabezado
\fancyhead[C]{%
    \begin{minipage}[c]{2.5cm}
        \centering
        \IfFileExists{imagenes/uta.png}{\includegraphics[width=2.2cm]{imagenes/uta.png}}{\textbf{UTA}}
    \end{minipage}%
    \hfill
    \begin{minipage}[c]{8cm}
        \centering
        \small {\textbf{UNIVERSIDAD TECNICA DE AMBATO}} \\[0.1cm]
        \scriptsize \textit{FACULTAD DE INGENIERIA EN SISTEMAS ELECTRONICA E INDUSTRIAL} \\[0.1cm]
        \scriptsize \textbf{CARRERA DE SOFTWARE} \\[0.1cm]
        \scriptsize \textbf{CICLO ACADEMICO: FEBRERO - JULIO 2026}
    \end{minipage}%
    \hfill
    \begin{minipage}[c]{2.5cm}
        \centering
        \IfFileExists{imagenes/fisei.png}{\includegraphics[width=2.2cm]{imagenes/fisei.png}}{\textbf{FISEI}}
    \end{minipage}%
}

\begin{document}
\begin{center}
    {\large \textbf{INFORME DE GUIA PRACTICA}} \\
\end{center}

\section{Portada}
\begin{flushleft}
\begin{tabular}{ l l }
    \textbf{Tema:} & Sistema BI de segmentacion KMeans y recomendacion hibrida \\
    \textbf{Unidad de Organizacion Curricular:} & Profesional. \\
    \textbf{Nivel y Paralelo:} & 6to - A \\
    \textbf{Alumnos participantes:} & Toala Steeven \\
    &Lopez Alexis\\
    \textbf{Asignatura:} & Inteligencia de Negocios \\
    \textbf{Docente:} & Mg. Ing. Ruben Nogales. \\
\end{tabular}
\end{flushleft}


\section{Informe}
\subsection{Objetivos}

\subsubsection{General}
Desarrollar un flujo integral de analitica de datos para alquiler de peliculas que permita cargar y enriquecer informacion en MongoDB, segmentar clientes con KMeans y generar recomendaciones explicables para soporte de decisiones.


\subsubsection{Especificos}
\begin{itemize}
    \item Integrar datos de hechos y dimensiones desde Excel hacia un modelo enriquecido y legible para negocio.
    \item Construir DataFrames agregados utiles para analisis de peliculas, categorias y perfiles de clientes.
    \item Implementar segmentacion de clientes usando KMeans con variables de actividad, gasto y variedad.
    \item Desarrollar un recomendador hibrido que combine filtrado colaborativo, contenido y ajuste por cluster.
    \item Mejorar la interpretabilidad mediante motivos de recomendacion y desglose del calculo del score final.
    \item Presentar resultados mediante interfaz grafica y visualizaciones de clusters orientadas a usuarios no tecnicos.
\end{itemize}

\subsection{Introduccion}
El proyecto implementa una arquitectura por capas para resolver un caso de negocio de recomendacion sobre datos de alquiler. Primero se construye un proceso ETL que toma \texttt{BI-FINAL.xlsx}, integra dimensiones (cliente, pelicula, categoria, tiempo, tienda, ciudad y actores) y publica un fact enriquecido en JSON y MongoDB. Como parte de la limpieza, se eliminan IDs tecnicos de la salida analitica para trabajar con referencias legibles.

Sobre esta base se generan DataFrames agregados desde MongoDB: resumen de peliculas, resumen de categorias y perfil de clientes. Estos artefactos alimentan dos componentes clave:
\begin{itemize}
    \item Segmentacion de clientes con KMeans en \texttt{logica/clusters.py}.
    \item Recomendacion hibrida en \texttt{logica/recomendacion.py}.
\end{itemize}

El recomendador combina once senales ponderadas, entre ellas similitud coseno cliente-cliente, afinidad por categoria, popularidad, ajuste temporal y ajuste por segmento KMeans. El score final se calcula como suma ponderada:

\[
\text{score\_final} = \sum_{i=1}^{11} w_i s_i
\]

Con la implementacion actual:

\begin{align*}
\text{score\_final} =
&\;0.30\,s_{colaborativo} + 0.16\,s_{categoria} + 0.11\,s_{popularidad} + 0.07\,s_{ingreso} \\
&+ 0.05\,s_{duracion} + 0.08\,s_{categoria\_global} + 0.07\,s_{ajuste\_cliente} \\
&+ 0.04\,s_{temporal} + 0.02\,s_{diversidad} + 0.06\,s_{cluster\_ajuste} + 0.04\,s_{cluster\_categoria}
\end{align*}

En la interfaz Tkinter, el usuario puede seleccionar cliente, generar recomendaciones, visualizar motivo principal y revisar el detalle matematico del score por pelicula. En paralelo, la vista de clusters entrega graficos en ventanas separadas con descripciones para facilitar interpretacion.

\subsection{Desarrollo tecnico detallado del sistema}

\subsubsection{1) Proceso ETL completo: desde Excel hasta JSON y MongoDB}
El flujo inicia en \texttt{logica/carga\_datos.py} con el archivo \texttt{BI-FINAL.xlsx}. Este modulo tiene como objetivo transformar una estructura transaccional con IDs tecnicos en un fact enriquecido, legible y listo para analitica.

\paragraph{Paso 1: lectura de hojas y control de existencia}
Se utiliza \texttt{pd.ExcelFile} para leer una sola vez el libro y luego extraer cada hoja requerida con la funcion \texttt{cargar\_hoja}. Si una hoja obligatoria no existe, el proceso se detiene con error explicito.

\begin{lstlisting}
if nombre_hoja not in excel.sheet_names:
    raise ValueError(f'No existe la hoja requerida: {nombre_hoja}')
\end{lstlisting}

\paragraph{Paso 2: limpieza tipologica del fact}
En \texttt{preparar\_fact} se convierten a numerico columnas de medida y claves (\texttt{ingreso}, \texttt{duracion\_alquiler}, \texttt{id\_*}) usando \texttt{errors='coerce'} para capturar inconsistencias sin romper el flujo.
Posteriormente se filtran registros sin \texttt{ingreso}, porque el sistema de recomendacion y los agregados dependen de este campo.

\paragraph{Paso 3: enriquecimiento dimensional}
Con una secuencia de \texttt{merge} por llave, el fact incorpora atributos descriptivos:
\begin{itemize}
    \item Cliente: nombre, apellido, email, estado activo.
    \item Pelicula: titulo, duracion, clasificacion, idioma, precio.
    \item Categoria, ciudad, tiempo y tienda.
\end{itemize}
Esto convierte datos orientados a almacenamiento en datos orientados a analisis y explicacion de negocio.

\paragraph{Paso 4: consolidacion de actores por pelicula}
Con la tabla puente \texttt{PUENTE\_PELICULA\_ACTOR} y \texttt{DIM\_ACTOR} se construyen:
\begin{itemize}
    \item \texttt{pelicula\_actores}: lista de actores unicos por pelicula.
    \item \texttt{cantidad\_actores}: total de actores unicos.
\end{itemize}
La agregacion usa \texttt{groupby('id\_pelicula')} para concentrar la informacion en un solo registro por pelicula.

\paragraph{Paso 5: identificador legible de cliente}
Se crea \texttt{cliente\_nombre\_completo} concatenando nombre y apellido, para evitar depender de IDs internos en reportes e interfaz.

\paragraph{Paso 6: salida final y persistencia}
Antes de guardar, se eliminan columnas que inician con \texttt{id\_}. El resultado se persiste en:
\begin{itemize}
    \item JSON: \texttt{fact\_alquiler\_proyecto.json}.
    \item MongoDB: \texttt{BI\_Final.fact\_alquiler}.
\end{itemize}
Se ejecuta \texttt{collection.drop()} para evitar duplicados entre corridas y garantizar consistencia de una carga completa por ejecucion.

\subsubsection{2) Como se construyen y para que sirven los DataFrames}
La capa \texttt{logica/dataframes.py} construye DataFrames agregados desde MongoDB para separar calculo analitico de la capa de interfaz.

\paragraph{DataFrame 1: \texttt{df\_resumen\_peliculas}}
Se agrupa por \texttt{pelicula\_titulo} y se calculan:
\begin{itemize}
    \item \texttt{total\_alquileres}: cuantas veces se alquilo.
    \item \texttt{ingreso\_promedio} y \texttt{ingreso\_total}: rentabilidad.
    \item \texttt{clientes\_unicos}: alcance real de audiencia.
    \item \texttt{duracion\_promedio}: comportamiento temporal del alquiler.
\end{itemize}
\textbf{Por que se usa:} es la tabla base para rankear candidatos en recomendacion, porque contiene popularidad, valor economico y comportamiento de consumo.

\paragraph{DataFrame 2: \texttt{df\_resumen\_categorias}}
Se agrupa por \texttt{categoria\_nombre} y se obtiene volumen, clientes e ingresos.
\textbf{Por que se usa:} permite medir fuerza global de categoria y evitar que una recomendacion dependa solo de similitud cliente-cliente.

\paragraph{DataFrame 3: \texttt{df\_perfil\_clientes}}
Se agrupa por \texttt{cliente\_ref} y se calcula:
\begin{itemize}
    \item \texttt{total\_alquileres} (actividad),
    \item \texttt{peliculas\_unicas} y \texttt{categorias\_unicas} (variedad),
    \item \texttt{gasto\_total} y \texttt{gasto\_promedio} (capacidad de gasto).
\end{itemize}
\textbf{Por que se usa:} es el vector de comportamiento que alimenta KMeans y tambien componentes del recomendador.

\paragraph{Normalizaciones aplicadas}
Para combinar metricas con escalas distintas se usa:
\[
\hat{x}=\frac{x-\min(x)}{\max(x)-\min(x)}
\]
Si \(\max(x)=\min(x)\), la serie se reemplaza por ceros para evitar division por cero.

En categorias se define un score combinado:
\[
\text{score\_categoria\_global}=0.6\cdot\text{popularidad\_categoria\_norm}+0.4\cdot\text{ingreso\_categoria\_norm}
\]

\subsubsection{3) Algoritmo de clusters: como se obtienen exactamente}
La segmentacion esta en \texttt{logica/clusters.py} y usa KMeans sobre \texttt{df\_perfil\_clientes}.

\paragraph{Variables de entrada al modelo}
\begin{itemize}
    \item \texttt{gasto\_promedio}
    \item \texttt{total\_alquileres}
    \item \texttt{peliculas\_unicas}
    \item \texttt{categorias\_unicas}
\end{itemize}

\paragraph{Pipeline del clustering}
\begin{enumerate}
    \item Conversion de columnas a numerico y eliminacion de filas invalidas.
    \item Estandarizacion con \texttt{StandardScaler} para que ninguna variable domine por escala.
    \item Definicion de \(k\) robusta: \texttt{k = min(n\_clusters, numero\_clientes\_validos)}.
    \item Entrenamiento de \texttt{KMeans(random\_state=42, n\_init=10)}.
    \item Etiquetado semantico de clusters por gasto promedio del centroide.
\end{enumerate}

\paragraph{Interpretacion de etiquetas}
Los clusters no se dejan como numeros tecnicos. Se renombran como:
\begin{itemize}
    \item \texttt{Mayor gasto}
    \item \texttt{Gasto medio}
    \item \texttt{Menor gasto}
\end{itemize}
segun el orden descendente del centro de \texttt{gasto\_promedio}.

\paragraph{Salida final de segmentacion}
Cada cliente queda con:
\texttt{cliente\_ref}, \texttt{cluster\_id}, \texttt{cluster\_nombre}, \texttt{cluster\_gasto\_centro}.
Esta salida se integra al recomendador y tambien se exporta en \texttt{clusters\_fact\_alquiler\_proyecto.csv}.

Fragmento representativo del codigo aplicado:
\begin{lstlisting}
scaler = StandardScaler()
X_scaled = scaler.fit_transform(base[columnas_modelo].fillna(0.0))

kmeans = KMeans(n_clusters=k, random_state=42, n_init=10)
base['cluster_id'] = kmeans.fit_predict(X_scaled)
\end{lstlisting}

\subsubsection{4) Algoritmo del sistema de recomendacion: desglose completo}
El modulo \texttt{logica/recomendacion.py} implementa un recomendador hibrido de 11 senales.

\paragraph{Etapa A: preparacion y matrices base}
\begin{itemize}
    \item Se construye \texttt{cliente\_ref} y \texttt{pelicula\_ref} en formato legible.
    \item Se arma matriz cliente-pelicula por conteo de alquileres.
    \item Se separan peliculas vistas y candidatas no vistas.
\end{itemize}

\paragraph{Etapa B: similitud colaborativa}
Se calcula similitud coseno cliente-cliente:
\[
\cos(\theta)=\frac{u\cdot v}{\|u\|\,\|v\|}
\]
Con los \texttt{n\_vecinos} mas similares se estima una preferencia ponderada por pelicula candidata.

\paragraph{Etapa C: senales de contenido y contexto}
Para cada candidato se calculan:
\begin{itemize}
    \item \texttt{score\_categoria}: afinidad del cliente con la categoria.
    \item \texttt{score\_categoria\_global}: fortaleza global de la categoria.
    \item \texttt{score\_popularidad}, \texttt{score\_ingreso}, \texttt{score\_duracion} desde resumen de peliculas.
    \item \texttt{score\_ajuste\_cliente}: cercania entre ingreso promedio de pelicula y gasto promedio del cliente.
    \item \texttt{score\_temporal}: similitud entre patron mensual del cliente y patron mensual de la pelicula.
    \item \texttt{score\_diversidad}: incentivo a categorias menos consumidas por ese cliente.
\end{itemize}

\paragraph{Etapa D: senales de cluster}
Usando la segmentacion KMeans:
\begin{itemize}
    \item \texttt{score\_cluster\_ajuste}: cercania entre ingreso de pelicula y centro economico del cluster.
    \item \texttt{score\_cluster\_categoria}: frecuencia de la categoria dentro del mismo cluster del cliente.
\end{itemize}

\paragraph{Etapa E: score final y ranking}
El ranking se obtiene por suma ponderada:
\begin{align*}
\text{score\_final}=
&\;0.30\,s_{colaborativo}+0.16\,s_{categoria}+0.11\,s_{popularidad}+0.07\,s_{ingreso} \\
&+0.05\,s_{duracion}+0.08\,s_{categoria\_global}+0.07\,s_{ajuste\_cliente} \\
&+0.04\,s_{temporal}+0.02\,s_{diversidad}+0.06\,s_{cluster\_ajuste}+0.04\,s_{cluster\_categoria}
\end{align*}

Fragmento representativo del codigo aplicado:
\begin{lstlisting}
candidatos['score_final'] = (
    candidatos['score_colaborativo'] * 0.30
    + candidatos['score_categoria'] * 0.16
    + candidatos['score_popularidad'] * 0.11
    + candidatos['score_ingreso'] * 0.07
    + candidatos['score_duracion'] * 0.05
    + candidatos['score_categoria_global'] * 0.08
    + candidatos['score_ajuste_cliente'] * 0.07
    + candidatos['score_temporal'] * 0.04
    + candidatos['score_diversidad'] * 0.02
    + candidatos['score_cluster_ajuste'] * 0.06
    + candidatos['score_cluster_categoria'] * 0.04
)
\end{lstlisting}

Luego se ordena de mayor a menor \texttt{score\_final} y se devuelven las \texttt{n\_recomendaciones} top.

\paragraph{Etapa F: explicabilidad}
Cada fila recomendada incorpora:
\begin{itemize}
    \item \texttt{motivo}: explicacion corta para lectura rapida.
    \item \texttt{motivo\_detalle}: desglose numerico tipo \texttt{valor x peso = aporte} por cada senal.
\end{itemize}
Esto permite auditar el por que de cada resultado.

\subsubsection{5) Como funciona la interfaz del sistema}
La vista \texttt{vistas/interfaz\_recomendacion.py} conecta usuario y motor de recomendacion.

\paragraph{Flujo de uso}
\begin{enumerate}
    \item Carga inicial desde MongoDB.
    \item Seleccion de cliente en \texttt{Combobox}.
    \item Definicion de numero de recomendaciones y vecinos.
    \item Ejecucion de \texttt{recomendar\_peliculas()}.
    \item Visualizacion en tabla (pelicula, categoria, score, motivo principal).
    \item Panel inferior con detalle matematico de la recomendacion seleccionada.
\end{enumerate}

\paragraph{Elementos de salida para el usuario}
\begin{itemize}
    \item \textbf{Resumen del cliente}: algoritmo, segmento, historico, vecino mas similar, categorias preferidas.
    \item \textbf{Tabla principal}: resultado compacto para decision rapida.
    \item \textbf{Detalle de score}: trazabilidad completa de aportes.
    \item \textbf{Exportacion CSV}: persistencia externa de recomendaciones.
\end{itemize}

\subsubsection{6) Visualizacion de clusters}
La descripcion puntual de cada grafico se presenta junto a su imagen correspondiente en la seccion de evidencia grafica, para facilitar la lectura directa de cada resultado.

\subsubsection{7) Integracion end-to-end del sistema}
El sistema opera con una trazabilidad clara:
\begin{enumerate}
    \item ETL integra y limpia Excel.
    \item Se genera JSON y se actualiza MongoDB.
    \item Desde Mongo se construyen DataFrames analiticos.
    \item Con perfil de clientes se ejecuta KMeans.
    \item El recomendador combina colaborativo + contenido + contexto + cluster.
    \item Interfaz y graficos exponen resultados con explicabilidad.
\end{enumerate}

\subsubsection{8) Robustez y controles incluidos}
Se incorporaron validaciones para mantener estabilidad:
\begin{itemize}
    \item validacion de hojas requeridas en origen,
    \item conversion segura de tipos con \texttt{errors='coerce'},
    \item manejo de dataset vacio en cada capa,
    \item proteccion ante division por cero en normalizaciones,
    \item fallback de identificadores legibles cuando faltan columnas tecnicas,
    \item mensajes de error orientados a diagnostico.
\end{itemize}

\subsubsection{9) Fragmentos de implementacion clave (codigo real)}
Como el evaluador puede no tener acceso al repositorio, se incluyen fragmentos centrales de la implementacion para respaldar tecnicamente el informe.

\paragraph{ETL: eliminacion de IDs tecnicos y persistencia dual}
\begin{lstlisting}
def guardar_json_y_mongo(df: pd.DataFrame) -> None:
    columnas_sin_ids = [
        col for col in df.columns
        if not col.lower().startswith('id_')
    ]
    df_salida = df[columnas_sin_ids]
    registros = df_salida.to_dict(orient='records')

    with open(RUTA_JSON, 'w', encoding='utf-8') as archivo:
        json.dump(registros, archivo, ensure_ascii=False, indent=4)

    client = MongoClient(MONGO_URI)
    db = client[DB_NAME]
    collection = db[COLLECTION_NAME]
    collection.drop()
    if registros:
        collection.insert_many(registros)
\end{lstlisting}

\paragraph{Carga inicial desde Excel y validacion de hojas requeridas}
\begin{lstlisting}
def cargar_hoja(excel: pd.ExcelFile, nombre_hoja: str) -> pd.DataFrame:
    if nombre_hoja not in excel.sheet_names:
        raise ValueError(f'No existe la hoja requerida: {nombre_hoja}')
    return excel.parse(nombre_hoja)


def leer_tablas() -> dict:
    excel = pd.ExcelFile(RUTA_EXCEL)
    tablas = {
        'fact': cargar_hoja(excel, 'FACT_ALQUILER'),
        'actor': cargar_hoja(excel, 'DIM_ACTOR'),
        'categoria': cargar_hoja(excel, 'DIM_CATEGORIA'),
        'ciudad': cargar_hoja(excel, 'DIM_CIUDAD'),
        'cliente': cargar_hoja(excel, 'DIM_CLIENTE'),
        'pelicula': cargar_hoja(excel, 'DIM_PELICULA'),
        'tiempo': cargar_hoja(excel, 'DIM_TIEMPO'),
        'tienda': cargar_hoja(excel, 'DIM_TIENDA'),
        'puente': cargar_hoja(excel, 'PUENTE_PELICULA_ACTOR'),
    }
    return tablas
\end{lstlisting}

\paragraph{Limpieza de campos numericos en la carga del fact}
\begin{lstlisting}
def preparar_fact(df: pd.DataFrame) -> pd.DataFrame:
    df = df.copy()
    columnas_numericas = [
        'id_hecho', 'id_tiempo', 'id_cliente', 'id_pelicula',
        'id_tienda', 'id_categoria', 'id_ciudad',
        'cantidad_alquiler', 'ingreso', 'duracion_alquiler'
    ]
    for columna in columnas_numericas:
        if columna in df.columns:
            df[columna] = pd.to_numeric(df[columna], errors='coerce')

    df = df.dropna(subset=['ingreso'])
    return df
\end{lstlisting}

\paragraph{Carga de datos para el recomendador desde MongoDB}
\begin{lstlisting}
def cargar_datos_mongo() -> pd.DataFrame:
    client = MongoClient(MONGO_URI)
    db = client[DB_NAME]
    collection = db[COLLECTION_NAME]
    df = pd.DataFrame(list(collection.find({}, {'_id': 0})))
    if df.empty:
        raise ValueError('No hay datos cargados en MongoDB para recomendar.')
    return df


def preparar_datos(df: pd.DataFrame) -> pd.DataFrame:
    df = df.copy()
    df['ingreso'] = pd.to_numeric(df['ingreso'], errors='coerce')

    if 'cliente_nombre_completo' in df.columns:
        df['cliente_ref'] = df['cliente_nombre_completo'].astype(str).str.strip()
    if 'pelicula_titulo' in df.columns:
        df['pelicula_ref'] = df['pelicula_titulo'].astype(str).str.strip()

    df = df.dropna(subset=['cliente_ref', 'pelicula_ref', 'ingreso'])
    return df
\end{lstlisting}

\paragraph{Agregacion en MongoDB para peliculas}
\begin{lstlisting}
resultado_2 = list(fact_alquiler.aggregate([
    {'$group': {
        '_id': '$pelicula_titulo',
        'pelicula_titulo': {'$first': '$pelicula_titulo'},
        'categoria_nombre': {'$first': '$categoria_nombre'},
        'total_alquileres': {'$sum': 1},
        'ingreso_promedio': {'$avg': '$ingreso'},
        'ingreso_total': {'$sum': '$ingreso'},
        'clientes_unicos': {
            '$addToSet': {
                '$ifNull': [
                    '$cliente_nombre_completo',
                    '$cliente_nombre'
                ]
            }
        }
    }},
    {'$project': {
        '_id': 0,
        'pelicula_ref': '$pelicula_titulo',
        'pelicula_titulo': '$pelicula_titulo',
        'categoria_nombre': {'$ifNull': ['$categoria_nombre', 'Sin categoria']},
        'total_alquileres': 1,
        'ingreso_promedio': 1,
        'ingreso_total': 1,
        'clientes_unicos': {'$size': '$clientes_unicos'}
    }}
]))
\end{lstlisting}

\paragraph{Construccion de matriz cliente-pelicula}
\begin{lstlisting}
def construir_matriz_clientes_peliculas(df: pd.DataFrame) -> pd.DataFrame:
    return pd.pivot_table(
        df,
        index='cliente_ref',
        columns='pelicula_ref',
        values='ingreso',
        aggfunc='count',
        fill_value=0,
    )
\end{lstlisting}

\paragraph{Similitud coseno cliente-cliente}
\begin{lstlisting}
def calcular_similitud_coseno(matriz: pd.DataFrame, cliente_id: str) -> pd.Series:
    vector_objetivo = matriz.loc[cliente_id].to_numpy(dtype=float)
    norma_objetivo = np.linalg.norm(vector_objetivo)
    similitudes = {}

    for otro_cliente, fila in matriz.iterrows():
        if otro_cliente == cliente_id:
            continue
        vector_otro = fila.to_numpy(dtype=float)
        norma_otro = np.linalg.norm(vector_otro)
        if norma_objetivo > 0 and norma_otro > 0:
            similitud = float(
                np.dot(vector_objetivo, vector_otro)
                / (norma_objetivo * norma_otro)
            )
            if similitud > 0:
                similitudes[str(otro_cliente)] = similitud

    return pd.Series(similitudes).sort_values(ascending=False)
\end{lstlisting}

\paragraph{Motivo detallado y trazabilidad del score}
\begin{lstlisting}
def construir_motivo_detallado(fila: pd.Series) -> str:
    factores = [
        ('Colaborativo (clientes similares)', 'score_colaborativo', 0.30),
        ('Afinidad por categoria', 'score_categoria', 0.16),
        ('Popularidad de la pelicula', 'score_popularidad', 0.11),
        ('Ingreso promedio de la pelicula', 'score_ingreso', 0.07),
        ('Duracion promedio', 'score_duracion', 0.05),
    ]
    lineas = ['Calculo del score final (suma ponderada):']
    for nombre, col, peso in factores:
        valor = float(fila.get(col, 0.0))
        aporte = valor * peso
        lineas.append(f"- {nombre}: {valor:.4f} x {peso:.2f} = {aporte:.4f}")
    return '\\n'.join(lineas)
\end{lstlisting}

\paragraph{Interfaz: carga de resultados en tabla}
\begin{lstlisting}
for idx, fila in recomendaciones.reset_index(drop=True).iterrows():
    self.tree.insert(
        '',
        tk.END,
        iid=str(idx),
        values=(
            fila.get('pelicula_titulo', ''),
            fila.get('categoria_nombre', ''),
            f"{float(fila.get('score_final', 0.0)):.4f}",
            self._resumir_motivo(fila.get('motivo', '')),
        ),
    )
\end{lstlisting}

\subsection{Metodologia aplicada}
El desarrollo se ejecuto de manera incremental, validando cada capa antes de avanzar a la siguiente para evitar errores acumulados.

\subsubsection{Fase 1: Integracion y calidad de datos}
\begin{itemize}
    \item Se definio un flujo ETL para convertir datos de origen en un fact enriquecido.
    \item Se normalizaron tipos numericos y se descartaron registros invalidos para evitar ruido en el modelado.
    \item Se reemplazaron identificadores tecnicos por referencias legibles para mejorar comprension en reportes y UI.
\end{itemize}

\subsubsection{Fase 2: Modelado analitico}
\begin{itemize}
    \item Se construyeron agregados de peliculas, categorias y clientes como capa semantica reusable.
    \item Se aplico KMeans sobre variables estandarizadas para evitar sesgo por escala.
    \item Se etiquetaron segmentos por gasto para hacer la segmentacion interpretable en lenguaje de negocio.
\end{itemize}

\subsubsection{Fase 3: Motor de recomendacion}
\begin{itemize}
    \item Se construyo matriz cliente-pelicula para similitud colaborativa.
    \item Se incorporaron senales de contenido y contexto temporal para reducir recomendaciones triviales.
    \item Se incluyo ajuste por cluster para alinear recomendaciones al perfil de segmento.
\end{itemize}

\subsubsection{Fase 4: Explicabilidad y visualizacion}
\begin{itemize}
    \item Se implemento motivo principal y desglose de aportes por componente del score.
    \item Se mejoraron textos y etiquetas de graficos para usuarios no tecnicos.
    \item Se separaron visualizaciones por ventana para facilitar lectura y comparacion.
\end{itemize}

\subsection{Arquitectura y flujo operacional}
\begin{enumerate}
    \item \textbf{Entrada:} archivo \texttt{BI-FINAL.xlsx} con tabla de hechos y dimensiones.
    \item \textbf{ETL:} enriquecimiento y persistencia en JSON + MongoDB.
    \item \textbf{Analitica base:} calculo de DataFrames agregados desde MongoDB.
    \item \textbf{Segmentacion:} agrupacion de clientes mediante KMeans.
    \item \textbf{Recomendacion:} ranking hibrido con score ponderado.
    \item \textbf{Presentacion:} interfaz Tkinter y graficos de interpretacion.
\end{enumerate}

\subsection{Validaciones realizadas}
\begin{itemize}
    \item Verificacion de carga correcta en MongoDB y consistencia del numero de registros.
    \item Validacion de columnas criticas para clustering y recomendacion.
    \item Revision de ejecutabilidad de scripts en \texttt{logica/} y \texttt{vistas/}.
    \item Comprobacion de explicaciones de score y legibilidad de interfaz.
\end{itemize}

\subsection{Resultados obtenidos}
\begin{itemize}
    \item Pipeline completo funcional: Excel $\rightarrow$ JSON/MongoDB $\rightarrow$ Clusters $\rightarrow$ Recomendaciones.
    \item Segmentacion estable y reusable para analisis y para mejorar el score del recomendador.
    \item Recomendaciones explicables, con trazabilidad de cada componente del score.
    \item Interfaz mejorada en legibilidad de motivos y detalle de calculo.
    \item Graficos de cluster reetiquetados para usuarios sin perfil tecnico.
\end{itemize}

\subsection{Evidencia grafica de resultados}
En esta seccion se dejan espacios para insertar capturas reales de la ejecucion del sistema.

\begin{figure}[H]
    \centering
    \espacioimagen{imagenes/graficoclientessegmento.png}{grafico de cantidad de clientes por segmento}
    \caption{Distribucion de clientes por segmento}
\end{figure}

\textbf{Descripcion:} Este grafico muestra el tamano relativo de cada cluster. La altura de cada barra representa la cantidad de clientes del segmento y permite identificar que grupos son predominantes y en cuales conviene priorizar estrategias comerciales.

\begin{figure}[H]
    \centering
    \espacioimagen{imagenes/graficosimilitudclientes.png}{grafico de similitud de comportamiento de clientes}
    \caption{Mapa de similitud de comportamiento de clientes}
\end{figure}

\textbf{Descripcion:} Esta visualizacion proyecta a los clientes en 2 dimensiones mediante PCA usando variables de actividad, gasto y variedad. Cada punto representa un cliente; puntos cercanos indican comportamientos similares. El tamano del punto esta asociado a \texttt{gasto\_promedio}.

\begin{figure}[H]
    \centering
    \espacioimagen{imagenes/grafico_perfil_segmento.png}{grafico comparativo de perfil por segmento}
    \caption{Comparacion de perfil por segmento}
\end{figure}

\textbf{Descripcion:} Este grafico compara el perfil medio de cada segmento en escala normalizada \([0,1]\): nivel de gasto, nivel de actividad, variedad de peliculas y variedad de categorias. Permite observar rapidamente fortalezas relativas de cada cluster sin sesgo por diferencias de escala.

\begin{figure}[H]
    \centering
    \espacioimagen{imagenes/grafico_resumen_segmento.png}{mapa de calor de resumen por segmento}
    \caption{Resumen rapido por segmento}
\end{figure}

\textbf{Descripcion:} El mapa de calor resume por segmento la cantidad de clientes, el gasto promedio y la actividad promedio en valores normalizados. Los colores mas intensos representan niveles mas altos y facilitan comparar el peso comercial de cada cluster en una sola vista.

\begin{figure}[H]
    \centering
    \espacioimagen{imagenes/interfaz_recomendacion.png}{captura de la interfaz de recomendacion}
    \caption{Interfaz de recomendacion con detalle de score}
\end{figure}

\textbf{Descripcion:} Este grafico muestra la interfaz del sistema de recomendacion con seleccion de cliente, tabla de peliculas recomendadas y panel de detalle explicativo. En ese panel se desglosa el \texttt{score\_final} por componentes, lo que permite interpretar por que una pelicula fue priorizada.

\subsection{Conclusiones}

\begin{itemize}
    \item Se logro una solucion BI integral, coherente y mantenible al separar responsabilidades entre modulos de logica y vistas.
    \item El uso de KMeans sobre perfiles de cliente aporta valor tanto analitico como operativo dentro del sistema de recomendacion.
    \item La explicabilidad del score final mejora la confianza en las recomendaciones y facilita su validacion por parte de usuarios de negocio.
    \item El ajuste de etiquetas y textos en visualizaciones fue clave para que la interpretacion no dependa de conocimientos tecnicos de programacion o modelado.
\end{itemize}


\subsection{Recomendaciones}

\begin{itemize}
    \item Incorporar seleccion automatica de numero de clusters (metodo codo/silhouette) para robustecer el modelado.
    \item Medir desempeno del recomendador con metricas offline como Precision@K y Recall@K.
    \item Agregar historico de ejecuciones para comparar evolucion de segmentos y cambios de comportamiento.
    \item Estandarizar la plantilla de reportes del proyecto para futuras iteraciones y evidencias academicas.
\end{itemize}


\subsection{Referencias}

\begin{thebibliography}{9}
\bibitem{sklearn}
scikit-learn Developers,
\textit{scikit-learn User Guide},
\url{https://scikit-learn.org/stable/user_guide.html}

\bibitem{pandas}
The pandas development team,
\textit{pandas Documentation},
\url{https://pandas.pydata.org/docs/}

\bibitem{pymongo}
MongoDB Inc.,
\textit{PyMongo Documentation},
\url{https://pymongo.readthedocs.io/}

\bibitem{matplotlib}
Matplotlib Development Team,
\textit{Matplotlib Documentation},
\url{https://matplotlib.org/stable/}
\end{thebibliography}



\vfill

\end{document}
